% !TeX program = tectonic
\documentclass[11pt,a4paper]{article}
\usepackage{fontspec}
\usepackage[english]{babel}
\babelfont{rm}[Language=Default]{Liberation Serif}
\babelfont[Language=Default]{sf}{Carlito}
\babelfont[Language=Default]{tt}{DejaVu Sans Mono}
\usepackage{amsmath,amssymb,amsthm}
\usepackage{geometry}
\geometry{a4paper, top=2.4cm, bottom=2.4cm, left=2.4cm, right=2.4cm}
\usepackage{booktabs,tabularx,enumitem,xcolor,tcolorbox}
\tcbuselibrary{breakable,skins}
\definecolor{linkblue}{HTML}{1F4E79}
\definecolor{statgreen}{HTML}{2F6B3A}
\definecolor{goldacc}{HTML}{8B7E5A}
\usepackage[colorlinks=true,linkcolor=linkblue,urlcolor=linkblue,unicode=true]{hyperref}
\theoremstyle{plain}
\newtheorem{theorem}{Theorem}
\newtheorem{lemma}[theorem]{Lemma}
\newtheorem{proposition}[theorem]{Proposition}
\newtheorem{corollary}[theorem]{Corollary}
\theoremstyle{definition}
\newtheorem{definition}[theorem]{Definition}
\theoremstyle{remark}
\newtheorem{remark}[theorem]{Remark}
\newtcolorbox{statusbox}[1][]{colback=statgreen!5,colframe=statgreen!75,
  title={\textbf{Summary of results:} #1},fonttitle=\small,boxrule=0.7pt,arc=2pt,breakable}
\newtcolorbox{databox}[1][]{colback=goldacc!6,colframe=goldacc!80,
  title={\textbf{Protocol data:} #1},fonttitle=\small,boxrule=0.7pt,arc=2pt,breakable}
\newtcolorbox{verifybox}[1][]{colback=linkblue!5,colframe=linkblue!80,
  title={\textbf{Verification:} #1},fonttitle=\small,boxrule=0.7pt,arc=2pt,breakable}
\widowpenalty=10000
\clubpenalty=10000
\setlength{\parskip}{0.35em}
\newcommand{\CC}{\mathbb{C}}
\newcommand{\RR}{\mathbb{R}}
\newcommand{\ZZ}{\mathbb{Z}}
\newcommand{\QQ}{\mathbb{Q}}
\newcommand{\Ga}{\Gamma}
\newcommand{\Bfun}{\mathrm{B}}
\newcommand{\im}{\operatorname{Im}}
\newcommand{\re}{\operatorname{Re}}
\newcommand{\lcm}{\operatorname{lcm}}
\newcommand{\SNF}{\operatorname{SNF}}
\newcommand{\Jac}{\operatorname{Jac}}
\newcommand{\rank}{\operatorname{rank}}
\newcommand{\Ind}{\operatorname{Index}}
\newcommand{\disc}{\operatorname{disc}}
\begin{document}
\begin{center}
{\LARGE\bfseries The Closed Form of the Period}\\[0.4em]
{\large the beta integral, the phase lattice, and equivariance}\\[0.6em]
{\normalsize Isaev Iskhak Khamzatovich}\\[0.2em]
{\normalsize The <<Dynamic Principle>> Program --- the \texttt{hodge-laboratory}}\\[0.15em]
{\small Standalone theorem monograph · Монография 2/16}
\end{center}
\vspace{0.4em}
{\color{linkblue}\hrule height 0.8pt}
\vspace{0.8em}

\section{Setting}

The periods of the Fermat curve --- integrals of holomorphic forms
over cycles --- admit a closed form: the transcendental part (a beta
integral) and the algebraic part (a root of unity) separate
explicitly. This separation is the technical foundation of the
whole $\Ga$-normalization.

\begin{lemma}[basis of holomorphic forms]\label{lem:forms}
The forms $\eta_{i,j}=\frac{x^i y^j dx}{y^{N-1}}$,
$i,j\ge0$, $i+j\le N-3$, are holomorphic on $C_N$; their number is
$g$, and they form a basis of $H^0(C_N,\Omega^1)$. Each
$\eta_{i,j}$ is an eigenform of the $\mu_N\times\mu_N$-action with
character $(a,b)=(i+1,j+1)$.
\end{lemma}

\begin{lemma}[closed form of the period]\label{lem:closedform}
In the coordinate $t=x^N$,
$\eta_{i,j}=\frac1N t^{\frac aN-1}(1-t)^{\frac bN-1}dt$. For a lift
$\gamma$ with winding data $(r,s)\in(\ZZ/N)^2$ the period is
\[
  P_{a,b}(r,s)=\frac1N\zeta_N^{ra+sb}\Omega_{a,b},
  \qquad
  \Omega_{a,b}=\frac{\Ga(\frac aN)\Ga(\frac bN)}{\Ga(\frac{a+b}N)}>0,
\]
and the image of the period functional lies in
$\frac{\Omega_{a,b}}{N}\ZZ[\zeta_N]$.
\end{lemma}

\begin{proof}
The substitution $t=x^N$ gives $dt=Nx^{N-1}dx$ and
$x^i\,dx=\tfrac1Nt^{\frac{i+1}N-1}dt$; and
$y^j/y^{N-1}=(1-t)^{\frac{j+1}N-1}$ since $y^N=1-t$. Each loop
around $t=0$ multiplies the branch $t^{1/N}$ by $\zeta_N$ --- the
integral by $\zeta_N^a$; a loop around $t=1$ gives $\zeta_N^b$. A
path with data $(r,s)$ carries $\zeta_N^{ra+sb}$ on the reference
integral $\frac1N\Bfun(a/N,b/N)$. Holomorphy: at $(\zeta,0)$,
$dx\sim-\zeta y^{N-1}dy$; at infinity in $u=X/Y$, $v=Z/Y$ the form
is $u^i(v\,du-u\,dv)v^{N-3-i-j}$, regular at $i+j\le N-3$.
Independence: forms of distinct characters are independent
(projection operators). Lattice: every cycle is an integer
combination of closed chains; the edge integral is a difference of
two closed-form values.
\end{proof}

\begin{corollary}[equivariance]
Translating a cycle by $(\zeta^u,\zeta^v)$ multiplies the period of
the eigenform $(a,b)$ by $\zeta_N^{ua+vb}$.
\end{corollary}

\begin{databox}[numbers]
$N=15$, $(a,b)=(1,1)$: $\Omega_{1,1}=246.897\ldots$;
$N=15$, $(7,8)$: $\Omega_{7,8}=\pi/\sin84^\circ=3.16298\ldots$ ---
a pure fractional share of $\pi$. Certificate B (V2): maximal
relative deviation from independent tanh-sinh over $69$ tests ---
$3.9\cdot10^{-36}$; phases (V3) --- deviation $0.0$.
\end{databox}

\begin{statusbox}[summary]
The closed form is proved by substitution and branch analysis;
equivariance is a direct corollary. The phase is fixed explicitly ---
no ``global phase'' is sought or postulated.
\end{statusbox}

\begin{verifybox}[reproduction]
\texttt{python3 laboratory.py} $\to$ ``Stand N=15/30 $\to$
V2/V3/V4''; the experiment designer accepts arbitrary $N$, $(a,b)$,
$(r,s)$ and prints the closed form, the independent integral, and
the verdict.
\end{verifybox}

\vfill
{\color{linkblue}\hrule height 0.6pt}
\vspace{0.5em}
{\small\color{gray}
License: individual exclusive license of Isaev Iskhak Khamzatovich.
All rights to the computations, formulas, and texts belong to the author.
Verification: \texttt{python3 laboratory.py} (``Stands''/``Certificates''),
Lean: \texttt{verification/lean/HodgeLaboratory.lean}.
}
\end{document}
